\section{Resultados e Análises}

\begin{frame}{Artefators Formalizados e Análise de Rastreabilidade}

    \begin{table}[h]
        \caption{Artefatos formalizados}
        \begin{center}
            \begin{tabular}{lc}
                \toprule
                \textbf{Categoria do Artefato} & \textbf{Quantidade} \\
                \midrule
                Casos de Uso (CdU) & 19 \\
                Regras de Negócio (RN) & 26 \\
                Dicionários de Dados (DD) & 4 \\
                \bottomrule
            \end{tabular}
        \end{center}
        
    {\footnotesize \textbf{Fonte:} Elaborado pelo autor.}
  \end{table}

    \begin{table}[h]
        \caption{Análise de rastreabilidade dos requisitos}
        \begin{center}
            \begin{tabular}{lccc}
                \toprule
                \textbf{Escopo Analisado} & \textbf{Total Mapeado} & \textbf{Sem Testes} & \textbf{Déficit} \\
                \midrule
                Casos de Uso & 19 & 15 & $\approx 79\%$ \\
                Fluxos Comportamentais & 79 & 65 & $\approx 82\%$ \\
                \bottomrule
            \end{tabular}
    
            {\footnotesize \textbf{Fonte:} Elaborado pelo autor.}
        \end{center}
    \end{table}
\end{frame}

\begin{frame}{Cobertura de Testes dos 4 CdUs}
    \begin{figure}[t]
    \centering 
    \label{fig:cobertura_fluxos}
    \caption{Cobertura de fluxos por Caso de Uso na Matriz de Rastreabilidade}
    \vspace{-1em}
    \resizebox{\textwidth}{!}{
    \begin{tikzpicture}
        \begin{axis}[
            xbar stacked, 
            width=14cm, 
            height=7cm, 
            enlarge y limits=0.15, % Dá um respiro maior nas margens superior e inferior
            xmin=0, 
            xlabel={Quantidade de Fluxos},
            ylabel={Casos de Uso},
            ylabel style={xshift=-10pt}, % Empurra o texto do eixo Y mais para a esquerda
            symbolic y coords={CdU10, CdU03, CdU02, CdU01},
            ytick=data, 
            % --- MELHORIAS ESTÉTICAS ---
            axis x line*=bottom, % Deixa apenas a linha do eixo X embaixo
            axis y line*=left,   % Deixa apenas a linha do eixo Y à esquerda
            xmajorgrids=true,    % Ativa a grade no eixo X
            grid style={dashed, gray!30}, % Deixa a grade tracejada e discreta
            bar width=18pt,      % Afina as barras para ficarem mais elegantes
            % Legenda sem borda e com mais espaço entre as palavras
            legend style={at={(0.5,-0.19)}, anchor=north, legend columns=-1, draw=none, /tikz/every even column/.append style={column sep=0.5cm}}, 
            % Código que imprime o número na barra apenas se ele for maior que zero
            nodes near coords={
                \pgfmathfloatifflags{\pgfplotspointmeta}{0}{}{\pgfmathprintnumber{\pgfplotspointmeta}}
            },
            nodes near coords align={center}, 
            % Formatação do número dentro da barra (negrito e tamanho menor)
            every node near coord/.append style={font=\footnotesize\bfseries, text=black!80},
        ]
        
        % 1ª SÉRIE: Fluxos 100% Cobertos (Verde com borda mais escura)
        \addplot[fill=mutedDGreen, draw=mutedDGreen!80!black, thick] coordinates {
            (2,CdU01) 
            (3,CdU02) 
            (1,CdU03)
            (1,CdU10)
        };
        
        % 2ª SÉRIE: Fluxos Parcialmente Cobertos (Laranja com borda mais escura)
        \addplot[fill=mutedOrange, draw=mutedOrange!80!black, thick] coordinates {
            (0,CdU01) % Este zero não será impresso, deixando o gráfico limpo!
            (3,CdU02) 
            (2,CdU03)
            (2,CdU10)
        };

        % 3ª SÉRIE: Fluxos Não Cobertos (Vermelho suave com borda)
        \addplot[fill=mutedRed, draw=mutedRed!80!black, thick] coordinates { 
            (2,CdU01) 
            (1,CdU02) 
            (1,CdU03)
            (4,CdU10)
        };
        
        \legend{Totalmente Cobertos, Parcialmente Cobertos, Não Cobertos}
        
        \end{axis}
    \end{tikzpicture}
    }
    {\footnotesize \textbf{Fonte:} Elaborado pelo autor.}
\end{figure}
\end{frame}

\begin{frame}{A Perspectiva Humana sobre a Dívida}
  \textbf{O cenário percebido:}
  \begin{itemize}
    \item A equipe desenvolvia e validava o SFM sem uma documentação formal.
    \item Não existe a elaboração de casos de teste sem que haja antes a implementação da funcionalidade, promovendo uma baixa cobertura dos requisitos elicitados.
    \item A equipe de \textit{design} assume o papel de analistas de requisitos.
  \end{itemize}

  \vspace{1.0em}

  \textbf{A Investigação de Percepção (O Questionário):}
  \begin{itemize}
    \item Como os desenvolvedores enxergavam a completude do sistema trabalhando sob essa Dívida Técnica?
    \item Qual o impacto real da introdução desses novos artefatos na rotina de trabalho?
  \end{itemize}
\end{frame}

\begin{frame}{Distribuição de Participantes nas Equipes do SFM}
    \input{graficos/pie_participation}
\end{frame}

\begin{frame}{Percepção Geral Sobre a Implementação das Funcionalidades}
    \input{graficos/percepcao_geral_funcionalidades} 
    %{\small \textbf{Fonte:} Elaborado pelo autor.}
\end{frame}

\begin{frame}{Percepção Sobre a Implementação do CdU 1}
    \input{graficos/cdu1}
    
    \textbf{Discussão:} O CdU 01 mostra que decisões baseadas puramente na observação visual induzem a equipe e a gestão ao erro. 
\end{frame}

\begin{frame}{Percepção Sobre a Implementação do CdU 3}
    \input{graficos/cdu3}
\end{frame}

\begin{frame}{Percepção Sobre a Implementação do CdU 8}
    \begin{figure}
    \centering
    \caption{Nível de implementação por fluxo do CdU 8}
    \vspace{-1em}
    \resizebox{\textwidth}{!}{
    \begin{tikzpicture}
        \footnotesize
        \begin{axis}[
            xbar stacked,
            width=12cm, 
            height=5cm, % Aumente se houver muitos fluxos
            xmin=0, xmax=7, % AJUSTADO: Agora travado em 7 respondentes
            xtick={0,1,2,3,4,5,6,7}, % AJUSTADO: Régua até 7
            xlabel={Número de Respondentes},
            ylabel={Fluxos do CdU 08},
            symbolic y coords={{Fluxo de Exceção 2}, {Fluxo de Exceção 1}, {Fluxo Alternativo 1}, {Fluxo Principal}},
            ytick=data,
            axis x line*=bottom,
            axis y line*=left,
            xmajorgrids=true,
            grid style={dashed, gray!30},
            bar width=18pt,
            legend style={
                at={(0.5,-0.22)}, 
                anchor=north, 
                legend columns=3, 
                draw=none,
                /tikz/every even column/.append style={column sep=0.5cm}
            },
            nodes near coords={
                \pgfmathfloatifflags{\pgfplotspointmeta}{0}{}{\pgfmathprintnumber{\pgfplotspointmeta}}
            },
            nodes near coords align={center},
            every node near coord/.append style={font=\footnotesize\bfseries, text=black!80},
        ]
        
        % Votos 1: Nada / Não Implementado (Vermelho)
        \addplot[fill=mutedRed, draw=mutedRed!80!black, thick] coordinates {
            (1,{Fluxo Principal}) 
            (1,{Fluxo Alternativo 1}) 
            (1,{Fluxo de Exceção 1})
            (1,{Fluxo de Exceção 2})
        };
        
        % Votos 2 (Laranja)
        \addplot[fill=mutedOrange, draw=mutedOrange!80!black, thick] coordinates {
            (2,{Fluxo Principal}) 
            (2,{Fluxo Alternativo 1}) 
            (1,{Fluxo de Exceção 1})
            (1,{Fluxo de Exceção 2})
        };
        
        % Votos 3 (Amarelo)
        \addplot[fill=mutedYellow, draw=mutedYellow!80!black, thick] coordinates {
            (0,{Fluxo Principal}) 
            (0,{Fluxo Alternativo 1}) 
            (1,{Fluxo de Exceção 1})
            (1,{Fluxo de Exceção 2})
        };
        
        % Votos 4 (Verde Claro)
        \addplot[fill=mutedLGreen, draw=mutedLGreen!80!black, thick] coordinates {
            (0,{Fluxo Principal}) 
            (0,{Fluxo Alternativo 1}) 
            (0,{Fluxo de Exceção 1})
            (0,{Fluxo de Exceção 2})
        };
        
        % Votos 5: Totalmente Implementado (Verde Médio com bom contraste)
        \addplot[fill=mutedDGreen, draw=mutedDGreen!80!black, thick] coordinates {
            (2,{Fluxo Principal}) 
            (1,{Fluxo Alternativo 1}) 
            (2,{Fluxo de Exceção 1})
            (2,{Fluxo de Exceção 2})
        };
        
        % Votos: DESCONHEÇO (Cinza Neutro)
        \addplot[fill=mutedGray, draw=mutedGray!80!black, thick] coordinates {
            (2,{Fluxo Principal}) 
            (3,{Fluxo Alternativo 1}) 
            (2,{Fluxo de Exceção 1})
            (2,{Fluxo de Exceção 2})
        };

        \legend{1 -- Nada Implementado, 2 -- Pouco Implementado, 3 -- Parcialmente Implementado, 4 -- Muito Implementado, 5 -- Totalmente Implementado, Desconheço Estado da Implementação}
        
        \end{axis}
    \end{tikzpicture}
    }
    
    {\footnotesize \textbf{Fonte:} Elaborado pelo autor.}
\end{figure}
\end{frame}

\begin{frame}{Percepção Sobre a Implementação do CdU 15}
    \input{graficos/cdu15}
\end{frame}

\begin{frame}{Discussão: O Gargalo na Qualidade}
  \textbf{A Dívida na Etapa de Testes:}
  \begin{itemize}
    \item A equipe de QA precisava aguardar a codificação completa para iniciar o planejamento e a escrita dos testes.
    \item Essa abordagem tardia e estritamente manual não é escalável e inviabiliza a futura automação.
  \end{itemize}

  \vspace{1em}
  
  \textbf{O Maior Risco Identificado:}
  \begin{itemize}
    \item A validação ocorria apenas quando a funcionalidade já estava totalmente construída.
    \item O perigo central não reside em falhas de código (\textit{bugs}), mas na possibilidade de implementar premissas que não correspondem às reais necessidades do projeto.
  \end{itemize}
\end{frame}

\begin{frame}{A Aceitação da Documentação}
    \begin{figure}[htpb]
    \centering
    \caption{Percepção da equipe sobre o grau em que a documentação contribuiria para a execução de suas tarefas diárias.}
    \vspace{-1em}
    \resizebox{\textwidth}{!}{
    \begin{tikzpicture}
        \footnotesize
        \begin{axis}[
            xbar stacked,
            width=12cm,
            height=6cm, 
            xmin=0, xmax=7, % Mantido para seus 7 respondentes
            xtick={0,1,2,3,4,5,6,7}, 
            xlabel={Número de Respondentes},
            ylabel={Casos de Uso Avaliados},
            symbolic y coords={{CdU 15 (Preliminar)}, {CdU 08 (Intermediário)}, {CdU 03 (Avançado)}, {CdU 01 (Avançado)}},
            ytick=data,
            axis x line*=bottom,
            axis y line*=left,
            xmajorgrids=true,
            grid style={dashed, gray!30},
            bar width=18pt,
            % Legenda centralizada para 5 itens
            legend style={
                at={(0.5,-0.22)}, 
                anchor=north, 
                xshift=-1.5cm,
                legend columns=3, % Ficará em 2 linhas
                draw=none,
                /tikz/every even column/.append style={column sep=0.5cm}
            },
            nodes near coords={
                \pgfmathfloatifflags{\pgfplotspointmeta}{0}{}{\pgfmathprintnumber{\pgfplotspointmeta}}
            },
            nodes near coords align={center},
            every node near coord/.append style={font=\footnotesize\bfseries, text=black!80},
        ]
        
        % Votos 1: Não Contribuiria (Vermelho)
        \addplot[fill=mutedRed, draw=mutedRed!80!black, thick] coordinates {
            (0,{CdU 01 (Avançado)}) 
            (0,{CdU 03 (Avançado)}) 
            (0,{CdU 08 (Intermediário)})
            (0,{CdU 15 (Preliminar)})
        };
        
        % Votos 2 (Laranja)
        \addplot[fill=mutedOrange, draw=mutedOrange!80!black, thick] coordinates {
            (1,{CdU 01 (Avançado)}) 
            (0,{CdU 03 (Avançado)}) 
            (0,{CdU 08 (Intermediário)})
            (0,{CdU 15 (Preliminar)})
        };
        
        % Votos 3 (Amarelo)
        \addplot[fill=mutedYellow, draw=mutedYellow!80!black, thick] coordinates {
            (0,{CdU 01 (Avançado)}) 
            (0,{CdU 03 (Avançado)}) 
            (1,{CdU 08 (Intermediário)})
            (0,{CdU 15 (Preliminar)})
        };
        
        % Votos 4 (Verde Claro)
        \addplot[fill=mutedLGreen, draw=mutedLGreen!80!black, thick] coordinates {
            (2,{CdU 01 (Avançado)}) 
            (1,{CdU 03 (Avançado)}) 
            (0,{CdU 08 (Intermediário)})
            (1,{CdU 15 (Preliminar)})
        };
        
        % Votos 5: Contribuiria Muito (Verde Escuro)
        \addplot[fill=mutedDGreen, draw=mutedDGreen!80!black, thick] coordinates {
            (4,{CdU 01 (Avançado)}) 
            (6,{CdU 03 (Avançado)}) 
            (6,{CdU 08 (Intermediário)})
            (6,{CdU 15 (Preliminar)})
        };
        
        % Alteramos a legenda para refletir a nova pergunta
        \legend{1 -- Não contribuiria, 2 -- Contribuiria Pouco, 3 -- Contribuiria Moderadamente, 4 -- Contribuiria Bastante, 5 -- Contribuiria muito}
        \end{axis}
    \end{tikzpicture}
    }

    {\footnotesize \textbf{Fonte:} Elaborado pelo autor.}
\end{figure}
\end{frame}

